\documentclass[10pt,a4paper]{amsart}
\usepackage[margin=19mm]{geometry}
\usepackage{amsmath,amssymb,mathtools}
\usepackage[scr=rsfs,cal=euler]{mathalfa}
\usepackage{xfrac}
\usepackage[unicode,hidelinks,bookmarks=false]{hyperref}
\newcommand{\ie}{i.e.\ }
\newcommand{\cf}{cf.\ }
\newcommand{\resp}{respectively\ }
\newcommand{\Subsection}{\subsection}
\providecommand{\nobreakdash}{\nobreak\hbox{-}}
\setlength{\emergencystretch}{2em}
\multlinegap0pt
\usepackage{bm}
\usepackage{bbm}
\usepackage{dsfont}
\usepackage{xspace}
\usepackage{cancel}
\usepackage{mathtools}
\usepackage{cleveref}
\usepackage{amsthm}
\makeatletter
\@ifundefined{thm}{\newtheorem{thm}{Thm}}{}
\@ifundefined{lem}{\newtheorem{lem}[thm]{Lemma}}{}
\@ifundefined{prop}{\newtheorem{prop}[thm]{Proposition}}{}
\makeatother
\title[A Choice-free look at algebraic extensions of valuations]{A Choice-free look at algebraic extensions of valuations}
\author{C\'edric A\"id}
\date{Source: arXiv:2511.07177v1 (CC BY 4.0)}
\begin{document}
\maketitle

\noindent\emph{Source and licence.} A Choice-free look at algebraic extensions of valuations. Version arXiv:2511.07177v1 (CC BY 4.0), distributed under the Creative Commons Attribution 4.0 International licence, as recorded in the arXiv licence metadata for this exact version and shown on the abstract page https://arxiv.org/abs/2511.07177v1. Full rights notice: licenses/arxiv-2511.07177-CC-BY-4.0.txt.\par\smallskip

\noindent\textit{Editorial scope.} Counted unit: the Prop whose source label is \texttt{loc} in the article (source0.tex lines 232--237, source label \texttt{loc}), delivered together with the article's own complete original proof of it (source0.tex lines 238--290, one \texttt{proof} environment of 53 lines). No mathematics is written, repaired or re-derived: statement and proof are the article's own text line for line. Required context retained uncounted: ax (lines 216--219). Every definition, display and label of those lines is retained unchanged. Omitted from this package and not redistributed: 1--215, 220--231, 291--715. Nothing inside a retained excerpt is omitted or altered: every retained body is delivered exactly as the article's lines are written, so no omission receipt of any kind accompanies this package. Citations: the retained excerpts carry no citation command at all, so no bibliography entry of the article is transcribed and no reference is redistributed. Reference adaptation: none was needed and none was made; every reference inside the delivered unit resolves inside it, so no unresolved reference or citation is produced. Printed slips kept literal: no printed word, symbol, hypothesis or number of the counted statement or of its proof was altered. Layout and class: the article's own class is \texttt{amsart} with packages including inputenc, amsmath, dsfont, amssymb, hyperref, cleveref, amsthm, biblatex; the wrapper is the lane's pinned \texttt{amsart} editorial wrapper at 10pt on A4, which restates the theorem-like environments the retained text needs and adds the article's own macro definitions that the retained text needs (none), with extra wrapper packages (bm, bbm, dsfont, xspace, cancel, mathtools, cleveref). The delivered numbering is the unit's own. Assets: no retained span contains a figure environment, a graphics-inclusion command, a PDF-page inclusion, a table or a TikZ picture; no image file is copied into this package, the package declares no asset dependency and no image file is redistributed with it. Licence: version arXiv:2511.07177v1 of this article is distributed under the Creative Commons Attribution 4.0 International licence, as recorded in the arXiv licence metadata for this exact version and shown on the abstract page https://arxiv.org/abs/2511.07177v1. A full rights notice is delivered at licenses/arxiv-2511.07177-CC-BY-4.0.txt. Characters: every retained excerpt is pure ASCII, so the delivered engine, XeLaTeX with its default Latin Modern text font, reports no missing character.
\begin{lem}{\label{ax}}
Let $R$ be a ring, $S$ an algebraic extension of $R$ and \hbox{$x\in S$.}
If $\sum_{i=0}^na_ix^i=0$ with $a_0,\ldots,a_n\in R$, then $a_nx$ is integral over $R$.
\end{lem}

\begin{prop}{\label{loc}}
Let $L$ be an algebraic extension of $K$ and $v$ a valuation on $K$.
Let $R$ be the relative integral closure of $\mathcal{O}_v$ in $L$.
If $I$ is a prime ideal of $R$ for which $\mathcal{O}_v\cap I=\mathfrak{m}_v$,
then $R_I$, the localization of $R$ at $I$, is a valuation ring of $L$ extending $\mathcal{O}_v$.
\end{prop}

\begin{proof}
Valuation rings of $L$ are characterized by the property that for every $x\in L^{\times}$, they contain at least one of $x$ or $x^{-1}$.
We will show that $R_I$ satisfies this property.

Let $x\in L^{\times}$.
We have that $\sum_{i=0}^na_ix^i=0$ for some $n\in\mathds{N}$ and $a_0,\ldots,a_n\in K$ with $a_0,a_n\neq0$.
Let $j$ be such that $$v(a_j)=\min\{v(a_i)\mid0\leq i\leq n\}.$$ Set $c_i=a_j^{-1}a_i$.
Then $\sum_{i=0}^nc_ix^i=0$ and $$\min\{v(c_i)\mid0\leq i\leq n\}=0,$$ by which $x$ is algebraic over $\mathcal{O}_v$.
Now set $$\alpha_j=\sum_{i=j}^{n}c_ix^{i-j}$$
for $0\leq j\leq n$. Then $$\alpha_{j-1}=x\alpha_j+c_{j-1}$$ and $$\alpha_{n}=c_{n}.$$
Set $\alpha_{n+1}=c_{n+1}=0$. Then $$\sum_{i=0}^{n+1}c_ix^i=0$$ and $$\alpha_{n}=x\alpha_{n+1}+c_{n}.$$
As $\alpha_{n+1}=0$, we have $x\alpha_{n+1}\in I$.
We will now proceed by reverse induction.

Let $1\leq j\leq n+1$ and suppose $x\alpha_j\in I$.
We can rewrite $$\sum_{i=0}^{n+1}c_ix^i=0$$
as
$$0=\sum_{i=0}^{j-1}c_ix^i+x^j\alpha_j=\sum_{i=0}^{j-2}c_ix^i+x^{j-1}(\alpha_jx+c_{j-1}).$$
Because $\alpha_jx\in I\subset R$ and $c_{j-1}\in R$, we get that $\alpha_{j-1}=\alpha_jx+c_{j-1}\in R$.
It follows from \Cref{ax} that $x\alpha_{j-1}$ is integral over $R$ and therefore also over $\mathcal{O}_v$.
As $R$ is the relative integral closure of $\mathcal{O}_v$ in $L$ and $x\alpha_{j-1}\in L$, we must also have that $x\alpha_{j-1}\in R$.
We now consider three cases:\par
\textit{Case $1$:} $v(c_{j-1})=0$\newline
Then $c_{j-1}\notin \mathfrak{m}_v=I\cap \mathcal{O}_v$.
Since $x\alpha_j\in I$ by hypothesis, $\alpha_{j-1}\notin I$ and so $\alpha_{j-1}$ is invertible in $R_I$.
We now have that $x=\frac{\alpha_{j-1}x}{\alpha_{j-1}}\in R_I$ and therefore we do not need to continue the induction.\par
\textit{Case $2$:} $x\alpha_{j-1}\notin I$\newline
Then $x\alpha_{j-1}$ is invertible in $R_I$.
We now have that $x^{-1}=\frac{\alpha_{j-1}}{\alpha_{j-1}x}\in R_I$ and therefore we do not need to continue the induction.\par
\textit{Case $3$:} $v(c_{j-1})>0$ and $x\alpha_{j-1}\in I$\newline
In this case, we see that the induction hypothesis applies to $j-1$ and we continue the induction.

As $\min\{v(c_i)\mid0\leq i\leq n\}=0$, case $1$ guarantees that the induction will stop.
When the induction stops we find that $x\in R_I$ or $x^{-1}\in R_I$.
This shows that $R_I$ is a valuation ring of $L$.

To show that $R_I$ is an extension of $\mathcal{O}_v$, we have to show that $R_I\cap K=\mathcal{O}_v$ and $IR_I\cap K=\mathfrak{m}_v$.
By construction, we have $\mathcal{O}_v\subseteq R\subseteq R_I$.
As we also have $\mathcal{O}_v\subseteq K$, this gives us the inclusion $\mathcal{O}_v\subseteq R_I\cap K$.
Next we show $\mathfrak{m}_v\subseteq IR_I\cap K$.
We have that
$$\mathfrak{m}_v=I\cap \mathcal{O}_v\subseteq I\subseteq IR_I,$$
which gives us the desired result.
Now we will show $R_I\cap K\subseteq \mathcal{O}_v$.
Let $x\in K\setminus \mathcal{O}_v$. Then $x^{-1}\in\mathfrak{m}_v\subseteq IR_I$.
The fact that $IR_I$ is a proper ideal of $R_I$ implies that $x\notin R_I$.
The last remaining inclusion to show is $IR_I\cap K\subseteq\mathfrak{m}_v$.
As we have already shown that $R_I\cap K=\mathcal{O}_v$, it suffices to show that $IR_I\cap \mathcal{O}_v\subseteq\mathfrak{m}_v$.
Let $x\in \mathcal{O}_v\setminus\mathfrak{m}_v$.
Then $x$ is invertible in $\mathcal{O}_v$.
This implies it is also invertible in $R_I$.
Therefore, $x\notin IR_I$.
\end{proof}

\end{document}
